% !TeX root = ../Report.tex
\chapter{System Engineering \& Requirement Analysis}

\section{Domain and Mission Context}

The ground segment of a space mission comprises all the facilities, software systems, communication infrastructure, and operational personnel responsible for supporting spacecraft operations throughout the mission lifecycle. It acts as the interface between the spacecraft and mission operators on Earth, enabling mission planning, spacecraft monitoring, command transmission, and anomaly management.

Within the Fucino Space Centre, the core ground segment is organized into several specialized subsystems, each responsible for a specific operational function. The main components include:

\begin{itemize}
    \item \textbf{Satellite Control Center (SCC)}, responsible for spacecraft monitoring and command execution;
    \item \textbf{Mission Planning Center (MPC)}, responsible for planning operational activities and scheduling mission tasks;
    \item \textbf{Telemetry, Tracking and Control (TT\&C)}, responsible for telemetry acquisition, spacecraft tracking, and command uplink;
    \item \textbf{Flight Dynamics}, responsible for orbit determination, maneuver planning, and trajectory analysis;
    \item \textbf{Mission Monitoring and Control (MMC)}, responsible for supervising mission execution and operational status;
    \item \textbf{Planning and Control Center (CPCM)}, including planning tools used to coordinate mission activities;
    \item \textbf{Communication Infrastructure (COMI)}, enabling data exchange between the different ground segment subsystems.
\end{itemize}

Although each subsystem performs a well-defined function, mission operations emerge from their continuous interaction. Operational data are generated by multiple heterogeneous sources and distributed across independent software systems. Consequently, obtaining a complete and coherent picture of the mission state requires correlating information originating from several subsystems operating simultaneously.

This distributed operational environment makes the ground segment inherently information-intensive and places significant cognitive demands on mission operators responsible for supervising mission execution.


\section{The Mission Manager's Role and Challenges}

Among the different operational roles, the Mission Manager is responsible for supervising the overall execution of the mission and coordinating the activities of the operational teams. The Mission Manager maintains situational awareness of the mission, evaluates operational events, coordinates responses to anomalies, and supports decision-making during both nominal and contingency operations.

To perform these tasks, the Mission Manager must continuously integrate information coming from multiple sources, including spacecraft telemetry, mission planning systems, flight dynamics analyses, communication systems, operational procedures, and reports generated by different operators. These information sources are often heterogeneous, asynchronous, and, in some situations, may provide incomplete or even conflicting evidence regarding the current mission state.

The increasing complexity of modern space missions has significantly amplified the amount of information that operators must process in real time. During critical operational phases or unexpected events, rapidly identifying relevant information and understanding its implications becomes particularly challenging. Maintaining situational awareness under these conditions requires considerable expertise and imposes a substantial cognitive workload on the Mission Manager.

These challenges motivate the adoption of intelligent Decision Support Systems capable of aggregating information from heterogeneous sources, assisting operators in assessing the current operational context, and supporting timely and informed decision-making.

The system we are going to develop is called FUNES: Foresight Understanding for Nominal and Emergency Scenarios. The objective of FUNES is to support mission managers during operational decision-making by reducing information fragmentation, improving situational awareness, and decreasing the cognitive workload associated with complex mission scenarios.

\section{FUNES System Overview}

FUNES (Foresight Understanding for Nominal and Emergency Scenarios) is an AI-based Decision Support Tool designed to support mission operations within an Earth Observation ground segment environment.

The system acts as an intermediate layer between existing mission operational systems and human operators, collecting heterogeneous operational data, analysing mission conditions, and providing decision-support outputs.

FUNES does not replace mission operators or execute autonomous spacecraft operations. Instead, it supports human decision-making by improving situational awareness, reducing information fragmentation, and providing recommendations based on available mission knowledge and operational data.

The main functions provided by FUNES include:

\begin{itemize}
    \item acquisition and processing of heterogeneous mission data;
    \item detection of anomalies and operational degradations;
    \item correlation of information from different mission domains;
    \item generation of operational recommendations;
    \item provision of explainable AI outputs to support operator decisions.
\end{itemize}

The system operates within the existing ground segment infrastructure and interfaces with external mission systems, including telemetry sources, mission planning systems, flight dynamics systems, monitoring tools, and operator interfaces.



\section{System and Software Requirement Specification according to ECSS}

Following the ECSS requirement engineering approach, the FUNES development process produced two main specification documents:

\begin{itemize}
    \item \textbf{System Requirements Document (SRD):} defining the technical requirements of the FUNES system as a complete entity within the mission ground segment environment;
    
    \item \textbf{Software Requirements Document (SRS):} defining the software-level requirements derived from the system specification and allocated to the software components.
\end{itemize}

The System Requirements Document represents the first level of technical specification. It describes the expected behaviour of the system independently from the implementation details and defines the capabilities required to satisfy operational needs.

The SRD includes:

\begin{itemize}
    \item \textbf{Functional requirements} describing the capabilities provided by the system;
    \item \textbf{performance requirements} defining measurable system constraints;
    \item \textbf{interface requirements} describing interactions with external and internal entities;
    \item \textbf{operational requirements} defining system behaviour during different operational modes;
    \item \textbf{safety, dependability, data management, and AI-specific requirements}.
\end{itemize}

According to ECSS principles, each requirement is written using a structured and verifiable formulation based on the "shall" statement. Requirements are required to be necessary, unambiguous, atomic, implementation-independent, verifiable, and traceable.

The Software Requirements Document is derived from the system specification by allocating system functions to software components. It defines the expected behaviour of the software architecture, including software functionalities, internal interfaces, data processing requirements, AI module interactions, and verification criteria.

The relationship between the two documents follows a hierarchical requirement flow where system-level requirements are progressively refined into software-level requirements while maintaining bidirectional traceability through a Requirement Traceability Matrix (RTM).

Citare ECSS-E-ST-10C Rev.1 per system engineering
ECSS-E-ST-40C Rev.1 per software engineering
ECSS-Q-ST-80C per software product assurance



\section{System Engineering Approach}

The development of the FUNES system followed a system engineering approach based on the principles defined by the ECSS.

The objective of the system engineering process was to transform operational needs identified within the mission control environment into a structured and traceable set of system and software requirements.

The adopted approach followed a top-down requirement derivation process, starting from the analysis of the operational scenario and stakeholder expectations, and progressively refining them into functional, performance, interface, operational, and AI-specific requirements.

The requirement definition process was performed according to the principles of ECSS-E-ST-10C, ensuring that each requirement was necessary, unambiguous, verifiable, traceable, atomic, and implementation-independent.

The identification of stakeholder needs represents the first step of the system engineering process. For the FUNES system, the primary stakeholders were identified within the mission operations environment, including Mission Managers, spacecraft operators, ground segment operators, and system maintenance personnel.

The analysis of the operational scenario highlighted several challenges associated with current mission operations. Mission operators must continuously monitor heterogeneous information sources, including spacecraft telemetry, mission planning information, ground segment status, flight dynamics data, network data, and operational reports.

During nominal and non-nominal operations, the integration and interpretation of such information require significant human effort, especially when rapid decisions are required. Therefore, the main operational need identified for FUNES was the capability to improve mission situational awareness by aggregating heterogeneous operational information and providing decision-support capabilities while preserving human authority over final operational decisions.



\section{System Requirements Definition}

The System Requirements Document (SRD) was developed to define the complete set of technical requirements governing the FUNES system behaviour and operational constraints.

The requirements were organized according to their functional role and operational impact. The main requirement categories identified were:

\begin{itemize}
    \item \textbf{Functional Requirements (FUN):} defining the capabilities provided by FUNES, including data acquisition, anomaly analysis, recommendation generation, explainability, and human-in-the-loop interaction;
    
    \item \textbf{Performance Requirements (PER):} defining measurable constraints related to accuracy, latency, reliability, scalability, and resource utilization;
    
    \item \textbf{Interface Requirements (INT):} defining interactions between FUNES and external mission systems, including telemetry sources, mission planning systems, databases, and operator interfaces;
    
    \item \textbf{Operational Requirements (OPE):} defining system behaviour during nominal, validation, training, and maintenance operations;
    
    \item \textbf{AI-specific Requirements (AI):} defining constraints related to AI lifecycle management, explainability, dataset governance, bias mitigation, and generative AI security.
\end{itemize}

All requirements were defined following ECSS-E-ST-10C principles, ensuring traceability between operational needs, system functions, and verification activities.



\section{Software Architecture} \label{sec:software architecture}

\section{Software Requirement Document}